\documentclass[10pt,a4paper]{amsart}
\usepackage[margin=19mm]{geometry}
\usepackage{amsmath,amssymb,mathtools}
\usepackage[scr=rsfs,cal=euler]{mathalfa}
\usepackage{xfrac}
\usepackage[unicode,hidelinks,bookmarks=false]{hyperref}
\newcommand{\ie}{i.e.\ }
\newcommand{\cf}{cf.\ }
\newcommand{\resp}{respectively\ }
\newcommand{\Subsection}{\subsection}
\providecommand{\nobreakdash}{\nobreak\hbox{-}}
\setlength{\emergencystretch}{2em}
\multlinegap0pt
\usepackage{amssymb}
\usepackage{mathtools}
\usepackage{float}
\usepackage{tikz}
\usepackage{algorithm}\usepackage{algpseudocode}
\newtheorem{theorem}{Theorem}
\newcommand{\PP}{\mathbb{P}}
\title{From Graph Laplacians to String Partition Functions: \\ A Rigorous Pathway from Discrete Spectra to Emergent Geometry}
\author{V.V. Tishkov}
\date{Source: arXiv:2605.00452v1 (CC BY 4.0)}
\begin{document}
\maketitle

\noindent\emph{Source and licence.} From Graph Laplacians to String Partition Functions: \\ A Rigorous Pathway from Discrete Spectra to Emergent Geometry. Version arXiv:2605.00452v1 (CC BY 4.0), distributed by arXiv under the Creative Commons Attribution 4.0 International licence, http://creativecommons.org/licenses/by/4.0/, as recorded in the arXiv licence metadata for this exact version and shown on the abstract page https://arxiv.org/abs/2605.00452v1. Full licence notice: \texttt{licenses/arxiv-2605.00452-CC-BY-4.0.txt}.\par\smallskip

\noindent\textit{Editorial scope.} Counted unit: the article's theorem thm:genus (source0.tex lines 107--113). The article's complete original proof of it is delivered with it (source0.tex lines 115--125, one \texttt{proof} environment of 11 lines). No mathematics is written, repaired or re-derived: the statement and the proof are the article's own lines, in the article's own notation, and no retained line is substituted or re-flowed.  No further source-local context is needed: every cross-reference of every retained line resolves inside the two delivered spans.  Nothing was omitted from any retained span, so the package carries no omission record.  No retained line contains a citation, so the wrapper carries no bibliography and no bibliography entry.  No retained line contains a figure environment, an image-inclusion command or drawing code, so no image is delivered and the package declares no asset dependency.  Editorial formatting only, in addition: an AMS \texttt{amsart} preamble that restates the article's own declarations and macros used by the retained lines, namely the environments \texttt{theorem} on separate counters, together with the standard packages those lines need.  The retained environments print in this document as Theorem 1 in order of appearance, whereas the article prints the same items with the numbering of its own full document.  The complete native source of the article as distributed in the arXiv e-print for this exact version is bundled unchanged as \texttt{sources/199563/source0.tex}.
\begin{theorem}[Genus of the Spectral Curve]\label{thm:genus}
The genus $g(X_G)$ of the spectral curve is given by
\[
g(X_G) = \left\lfloor \frac{d-1}{2} \right\rfloor,
\]
where $d$ is the number of distinct eigenvalues of $L_G$.
\end{theorem}

\begin{proof}
The projection $\pi: X_G \to \PP^1$ given by $(z,w) \mapsto z$ is a double cover branched precisely at the points $z = \mu_k$ for $k = 1,\dots,d$.

If $d$ is even, the point at infinity is not a branch point. By the Riemann-Hurwitz formula,
$2g(X_G) - 2 = 2(0 - 2) + d$ implies $g(X_G) = \frac{d-2}{2}$.

If $d$ is odd, the point at infinity is also a branch point, contributing an additional $1$ to the branching divisor. Then
$2g(X_G) - 2 = 2(0 - 2) + (d + 1)$ implies $g(X_G) = \frac{d-1}{2}$.

The floor function $\lfloor (d-1)/2 \rfloor$ unifies both cases.
\end{proof}

\end{document}
